\scanpage{30}
% archive sheet number: 29
\deleted{qui décrivent $\unsure{\mathfrak S_3|\mathscr S_3|G_3}\times\mathbb Z/2$, on va relever ces générateurs en des éléments de $\unsure{\mathscr E|\mathcal E}$, de façon que $l_\infty,l_1,l_0$ s'expriment comme des mots en ces éléments, et il \unsure{suffira|suffit} alors d'ajouter cette relation $l_\infty l_1l_0$ aux relations correspondant aux précédentes\unsure{....}}

$\dot\sigma_0,\dot\sigma_1,\dot\sigma_\infty$, $\unsure{\dot\rho|\rho},\tau$, avec les relations
\[
\left\{\begin{aligned}
\dot\sigma_0^2&=\dot\sigma_1^2=\dot\sigma_\infty^2=1\\
\dot\sigma_0\dot\sigma_1&=\dot\sigma_1\dot\sigma_\infty=\dot\sigma_\infty\dot\sigma_0=\dot\rho\\
\dot\rho^3&=1\\
\tau^2&=1\\
\tau&\text{ commute à }\dot\sigma_0,\dot\sigma_1,\dot\sigma_\infty\quad(\text{donc aussi à }\dot\rho)
\end{aligned}\right.
\]
on les relève en $\deleted{\tau_0,\tau_1,\tau_\infty}\sigma_0,\sigma_1,\sigma_\infty,\rho,\underbrace{\tau_0,\tau_1,\tau_\infty}_{\text{sur }\tau}$ avec les relations
\[
\left\{\begin{aligned}
\sigma_0^2&=\sigma_1^2=\sigma_\infty^2=\tau_0^2=[\tau_1^2=\tau_\infty^2]=1\\
\rho^3&=1\\
\rho\tau_0\rho^{-1}&=\tau_1,\quad\rho\tau_1\rho^{-1}=\tau_\infty,\quad\rho\tau_\infty\rho^{-1}=\tau_0\\
\tau_0\sigma_0&=\sigma_0\tau_0,\quad\tau_1\sigma_1=\sigma_1\tau_1,\quad\tau_\infty\sigma_\infty=\sigma_\infty\tau_\infty\\
\sigma_0\sigma_1&=\tau_\infty\tau_1\rho,\quad\sigma_1\sigma_\infty=\tau_0\tau_\infty\rho,\quad\sigma_\infty\sigma_0=\tau_1\tau_0\rho
\end{aligned}\right.
\]
% slanted
\margin{\deleted{\unsure{NB}} \unsure{Il faut voir si|Il faut vérifier si} les relations $\left\{\begin{aligned}\rho\sigma_0\rho^{-1}&=\sigma_1\\\rho\sigma_1\rho^{-1}&=\sigma_\infty\\\rho\sigma_\infty\rho^{-1}&=\sigma_0\end{aligned}\right.$ sont conséquences des relations ci-contre.}
\margin{Ce sont les éléments $\widetilde\sigma_0,\widetilde\sigma_1,\widetilde\sigma_\infty$ $\longrightarrow$}
\margin{$\bigl[$ NB \unsure{On trouve} $\begin{aligned}l_0&=\tau_\infty\tau_1,\quad\deleted{l_1=\tau_0\tau_\infty,}\\l_1&=\tau_0\tau_\infty\deleted{\unsure{l_\infty=}}\\l_\infty&=\tau_1\tau_0\end{aligned}$ et $l_\infty l_1l_0=1$ est automatique $\bigr]$}
\margin{\deleted{\unsure{NB}} Ceci implique que $\tau_0,\tau_1,\tau_\infty$ \unsure{conjuguent|engendrent} \unsure{le|un} ss-groupe engendré par $l_0,l_1,l_\infty$ $\uparrow$}

\deleted{\uline{Question} \uline{\deleted{Peut-on}\inserted{On peut}} exprimer $\rho$ en fonction de $\sigma_0,\sigma_1,\sigma_\infty,\tau_0,\tau_1,\tau_\infty$ ??, \deleted{d'après les} \inserted{l'une des} derniers groupes de relations, et substituer dans les autres.}
